\section{Глава~\ref{sec:synthesis}. Вся машина}

Глава-сводка новых опытов не приводит: каждая её строка опирается на главу, где утверждение разобрано, и на те же источники; твёрдо установлены шина \nt{ГЛУТ} и управление потоком через \nt{ГАМК}, а роли широковещательных каналов и их совместная работа --- в основном гипотеза.

{\footnotesize
\setlength{\tabcolsep}{3pt}\renewcommand{\arraystretch}{1.15}
\begin{longtable}{>{\raggedright}p{0.5cm}>{\raggedright}p{4.0cm}>{\raggedright}p{5.6cm}>{\raggedright}p{2.3cm}>{\raggedright\arraybackslash}p{1.7cm}}
\toprule
№ & Утверждение & Опыт и что измерено & Источник & Статус \\
\midrule\endfirsthead
\toprule
№ & Утверждение & Опыт и что измерено & Источник & Статус \\
\midrule\endhead
1 & На $8.6\cdot10^{10}$ нейронов приходится всего четыре рода управляющих сигналов~\eqref{eq:machine} & Подсчёт ядер клеток в гомогенате мозга (изотропный фракционатор): у взрослого мужчины $86.1\pm8.1$ млрд нейронов & \cite{Azevedo2009} & измерено \\
2 & Утверждение~\ref{prop:architecture}, первые два слоя: данные идут по адресной шине \nt{ГЛУТ}, знак связи задаёт отправитель, поток направляют и нормируют \nt{ГАМК}-нейроны (главы~\ref{sec:neurontypes}, \ref{sec:gaba}) & Один главный медиатор на нейрон (обзор измерений); прямое торможение сужает окно, в котором пирамидный нейрон складывает входы; деление на общую активность измерено во многих областях & \cite{Strata1999,Pouille2001,Carandini2012} & измерено \\
3 & Элементы ненадёжны: синапс теряет большую часть импульсов (табл.~\ref{tab:computer}, глава~\ref{sec:code}) & Срезы гиппокампа, минимальная стимуляция: больше половины импульсов не вызывают ответа в CA1; сбои порога и проведения по аксону исключены, причина --- вероятностный выброс медиатора & \cite{AllenStevens1994} & измерено \\
4 & Мозг работает на $\sim20$ ваттах, в среднем около импульса в секунду на нейрон (глава~\ref{sec:code}); такой машины инженеры пока не построили & Оценка~\eqref{eq:energy} из измеренных затрат: около $10^9$ молекул АТФ на импульс вместе с работой синапсов (кора грызунов); подробная оценка для коры даёт ещё меньшую частоту. Состояние техники --- по обзору & \cite{Attwell2001,Lennie2003,MehonicKenyon2022} & теория \\
5 & Обучение большинства синапсов --- произведение местного следа и одного общего числа $\delta$ (второе уравнение~\eqref{eq:machine}, глава~\ref{sec:dofa}) & \nt{ДОФА}-нейроны обезьяны отвечают на ошибку предсказания награды; в срезах полосатого тела дофамин увеличивает шипик, только если приходит через $0.3$--$2$\,с после глутаматного входа --- след измерен & \cite{Schultz1997,Yagishita2014} & измерено \\
6 & Отдельный провод ошибки к каждому выходу есть лишь в цепи обучения движениям (глава~\ref{sec:motorlearning}) & Мозжечок: совпадение сигнала лазящего волокна со входом ослабляет этот вход; у обезьяны сложные импульсы клеток Пуркинье несут направление ошибки движения и вызывают поправку & \cite{Ito2001,Herzfeld2018} & измерено \\
7 & Каждый широковещательный канал --- жгут из нескольких линий (утверждение~\ref{prop:bundle}) & Для \nt{ДОФА}: в одной задаче разные нейроны несут награду, движение и признаки стимула; быстрая ошибка и медленный уровень дофамина различны. Для \nt{СЕРО}, \nt{НОРА}, \nt{АЦЕТ} такая раскладка изучена хуже & \cite{Engelhard2019,Hamid2016,Mohebi2019} & измерено \\
8 & \nt{СЕРО}: регулятор уровня и, по гипотезе, горизонт $\tau_\gamma$ (глава~\ref{sec:serocomputer}) & Самоторможение: агонисты собственных рецепторов 5-HT$_{1A}$ тормозят серотониновые нейроны шва --- измерено. Горизонт предложен в теории; сильнейший опыт --- оптогенетическая активация этих нейронов у мыши удлиняет ожидание отложенной награды & \cite{Sprouse1987,Doya2002,Miyazaki2014} & гипотеза \\
9 & \nt{НОРА}: усиление $g$ выбирает между поиском и решением (глава~\ref{sec:nora}) & Рост отношения сигнала к шуму: у бодрствующей обезьяны норадреналин в слуховой коре подавляет фоновую активность сильнее, чем ответ на голоса сородичей --- измерено; мультипликативное усиление --- модель; роль в выборе между поиском и решением --- теория & \cite{Foote1975NE,ServanSchreiber1990,AstonJones2005} & гипотеза \\
10 & \nt{АЦЕТ}: переключает запись и чтение (глава~\ref{sec:acet}) & Три действия (ослабление внутренних связей, усиление входа, облегчение пластичности) измерены; переключение режимов косвенно поддержано опытом со скополамином у человека & \cite{HasselmoBower1992,Gil1997,Atri2004} & гипотеза \\
11 & Формула~\eqref{eq:machine} описывает машину целиком & Это сводка моделей глав~\ref{sec:glutcomputer}--\ref{sec:acet}, каждая часть опирается на свою главу; как целое опытом не проверялась & вывод в главе & гипотеза \\
12 & Ручки работают согласованно без общего управления: ошибка, неожиданность, запись, возврат (рис.~\ref{fig:timeline}) & Опыта нет: схема качественная. Отдельные звенья --- теория разделения труда \nt{НОРА} и \nt{АЦЕТ} и связь зрачка со скоростью обучения у человека & \cite{YuDayan2005,Nassar2012} & гипотеза \\
13 & Обучение по одному общему сигналу медленнее, чем больше нейронов влияют на результат (утверждение~\ref{prop:broadcastcost}) & Расчёт кривых обучения линейной сети: обучение по возмущению выходов медленнее градиентного спуска во столько раз, сколько выходов & \cite{Werfel2005} & теория \\
14 & Как кора настраивает многослойные цепи, неизвестно; кандидаты --- пачки импульсов, вычисления в ветвях нейрона, самообучение предсказанием & Пачки как носитель ошибки --- модель; повторить ответ подробной модели нейрона коры смогла лишь сеть в $5$--$8$ слоёв --- модель; кальциевые импульсы в ветвях нейронов коры человека, дающие исключающее ИЛИ, --- измерено в срезах; самообучение --- обзор свидетельств & \cite{Payeur2021,Beniaguev2021,Gidon2020,Keller2018} & гипотеза \\
\bottomrule
\end{longtable}}
